\documentclass[leqno, a4paper, 12pt]{jsarticle}
\usepackage{amssymb}
\usepackage{amsthm}
\usepackage{amsmath}
\usepackage{comment}
\usepackage{url}

\usepackage{enumitem}
%\usepackage{showkeys}
%定理環境 thmはあまり使わずpropをなるべく使う。
%主張は基本的に証明の中で使い、そのため番号を付けない。
%注意環境を付けてみた。
\newtheorem{thm}{定理}
\newtheorem{prop}[thm]{命題}
\newtheorem{cor}[thm]{系}
\newtheorem{lem}[thm]{補題}
\newtheorem{df}[thm]{定義}
\newtheorem{exam}[thm]{例}
\newtheorem{fact}[thm]{事実}
\newtheorem*{claim}{主張}
\newtheorem*{cau}{注意}
\newtheorem*{rmk}{注意}
\renewcommand{\proofname}{証明}
%矢印たち。
%\toは\rightarrowと同じ。\getsは\leftarrowと同じ。同値は\iff
\newcommand{\To}{\Longrightarrow}
\newcommand{\Gets}{\LongLeftarrow}
%黒板太字
\newcommand{\nn}{\mathbb{N}}
\newcommand{\zz}{\mathbb{Z}}
\newcommand{\qq}{\mathbb{Q}}
\newcommand{\rr}{\mathbb{R}}
\newcommand{\cc}{\mathbb{C}}
%無限大
\newcommand{\mgn}{\infty}
%差集合は\setminusで統一する。等号の否定は\neq
%真部分集合は\subsetneqq　偏微分は\partial
%ディスプレイスタイル。\dfracでディスプレイ分数
\newcommand{\dpst}{\displaystyle}

\title{素数の無限性}
\author{ナンブ　キトラ}
\date{\today}
\begin{document}
\maketitle
この文章では素数の無限性を証明する．
\begin{df}
$n, m\in \zz_{\ge 1}$とする．
$n=km$となる$k\in \zz_{\ge 1}$が存在するときに$m$は$n$の\textbf{約数}という．
$p\in \zz_{\ge 1}$が\textbf{素数}であるとは$p\neq 1$であり尚且つ$p$が$1$と自分自身以外の正の約数を持たないときにいう．\textbf{$n$の素因数}とは素数であり$n$を割り切れる数のことである．そして
$n\in \zz_{\ge 1}$に対して$1<a, b<n$を満たす
$a, b\in \zz_{\ge 1}$が存在して$n=ab$を満たすならば
$n$を\textbf{合成数}と呼ぶ．
\end{df}
合成数の定義から次がわかる．
\begin{prop}\label{nab}
$n$が合成数であることと$n\neq 1$で尚且つ$n$が素数でないことは同値である．
\end{prop}

\begin{prop}\label{soinsuu}
$n\in \zz_{\ge 1}$が$n\neq 1$を満たすならば$n$は素因数を持つ．
\end{prop}
\begin{proof}
まず
$A =\{\,  n\in \mathbb{N} 
\mid 
\text{
 $n\neq1$かつnは素因数を持たない}
 \, \}$
とおく．
命題は$A= \emptyset$と同値である.
今
$A\neq \emptyset$と仮定すると$A$の最小元が存在する．$a=\min A$と置く．
$a\in A$なので$a\neq 1$であり，$a$は素数ではない．よって命題 \ref{nab}より$a$は合成数である．
ゆえに$a=xy\ (1<x,y<a)$と書ける．
$a$の最小性から$x,y\not\in A$である．よって$x,y$は素因数をつことから$a$も素因数を持つことになり矛盾する．
\end{proof}


\begin{thm}
素数は無限に存在する．
\end{thm}
\begin{proof}
任意の$n>1$に対してそれよりも大きい素数が存在することを言えば良い．
$n!$を$1$から$n$までの間の自然数の積とする．このとき
$n!+1$は命題\ref{soinsuu}より素因数を持つ．
その素因数の一つを$p$とする．
$(n!+1)-n!=1$なので，
もしも$p$が$n!$の約数だとすると$p$は$1$を割り切ることになり
$p=1$となる．これは$p$が素数であることに矛盾する．よって$p$は$n!$の約数ではないので，$n!$の定義から$n<p$である.
\end{proof}











\end{document}